\section{Hoja de ruta de implementación en ingeniería: del prototipo de investigación al sistema en producción}

\subsection{Etapa uno: plataforma de experimentación reproducible}
Primero se construye un entorno reproducible, un protocolo de acciones unificado y scripts de evaluación estandarizados. El objetivo de esta etapa no es perseguir la puntuación más alta, sino establecer una base de experimentación diagnosticable, regresionable y escalable.

\subsection{Etapa dos: ciclo de retroalimentación de datos}
Con las tareas fallidas reales como núcleo, se retroalimentan de manera continua a la construcción de datos y al proceso de entrenamiento. El impulso guiado por tutoriales y los datos sintéticos avanzan en paralelo, pero siempre manteniendo la restricción de verificación de ejecución, para evitar el desplazamiento de la distribución de los datos.

\subsection{Etapa tres: vigilancia en producción}
Se incorporan control de permisos, auditoría de operaciones, reversión ante anomalías y mecanismos de intervención manual. Solo cuando se cuenta con «observabilidad + capacidad de recuperación» puede el agente pasar de la demostración a la cadena crítica del negocio.

\begin{knowledgebox}{Requisitos de colaboración a nivel organizacional}
Un agente multimodal no es un proyecto que un solo equipo de modelos pueda completar por sí solo. Requiere la colaboración de ingeniería de entornos, ingeniería de evaluación, ingeniería de datos, entrenamiento de modelos y equipos de seguridad de producto, para formar un ritmo de lanzamiento unificado.
\end{knowledgebox}

\begin{importantbox}{Principio de implementación en una frase}
\textit{``Primero hay que ver con claridad el proceso de las fallas, y después buscar puntuaciones más altas.''} Esto aumenta la productividad real más que amontonar parámetros o muestras a ciegas.
\end{importantbox}

\subsection{Resumen de la sección}
La ruta de implementación debe seguir el principio de ``la plataforma primero, el ciclo de retroalimentación de datos, la vigilancia en producción''. El valor central del curso está en ofrecer una ruta de ingeniería de agentes ejecutable, escalable y auditable.

